\documentclass{article}
\usepackage[margin=1in]{geometry}
\usepackage{amsmath}
\usepackage{enumitem}
\begin{document}
\title{Response to Reviewers}
\author{Abdelrahman Ellithy}
\date{December 02, 2025}
\maketitle
Dear Editor,
Thank you for the opportunity to resubmit our manuscript. We have carefully revised the paper according to the reviewers’ comments. Below is a point-by-point response:

\textbf{Reviewer 1}
\begin{enumerate}
\item This is a summary point. Specific issues are addressed below. Overall, the manuscript now includes theoretical analysis (Section IV, pages 11-13), adaptive parameters (pages 4, 13), and detailed benchmarks (pages 15-30), improving quality. However, some gaps like limited statistical depth persist.
\item Novelty emphasized in Related Work (Section II, page 2) as a "deterministic multi-phase architecture" differing from conditional hybrids like Brent. Added Section IV (pages 11-13) with convergence order analysis, theorem on local order (1 < p < 2), error bounds, and graceful degradation. Includes math like error relations and Taylor expansions. This goes beyond basic combinations.
\item Experimental setup (Section V, pages 13-14) now uses SymPy for symbolic definition but lambdifies to NumPy backend for numerical execution, with pre-compilation to reduce overhead. Timing uses \texttt{time.perf\_counter\_ns()} over 10,000 repetitions with 100 iteration warm-up and Garbage Collector (GC) disabled for reliability. CPU times are reported as secondry metric (e.g., 10.48 $\mu$s average, pages 1, 25), but shifted the focus to Number of Function Evaluations (NFE) (pages 1, 22, 26) as an implementation-independent main metric.
\item Benchmark now on 48 problems (pages 15-30), comparing all methods in same Python/NumPy environment (page 14). Reports NFE, iterations, FLOPs, and wins (pages 21-26). Statistical rigor added: averages over 10,000 runs per problem (page 14) wth other Statistical metrics like min, max, sd and median.
\item Formulas corrected, all algorithms follow standard format (e.g., trisection on page 3: $x_1 = a + (b-a)/3$, $x_2 = b - (b-a)/3$). Notation consistent (e.g., $\delta$ adaptive formula on pages 4, 12; $\epsilon=10^{-15}$ on page 3). Pseudocode vars standardized (e.g., fxB, fxFP, fxS across Algorithms 1-4, pages 5-10). No "ffp" or obvious irregularities found.
\item Pseudocode now uniform: All algorithms return iteration count \(n\), root \(x\), \(f(x)\), final \([a,b]\) or None (pages 5--10). Bracketing is maintained via sign checks (e.g., if \(f_a \cdot f_{xFP} < 0\), update \(b = x_{FP}\)). The modified secant is accepted only if \(a < x_S < b\) and \(|f_{xS}| < |f_{xFP}|\) (page 7). Termination criteria: \(|f(x)| \leq \text{tol}\) or max\_iter reached, with a midpoint fallback.
\item Benchmark includes applied examples: increased the test dataset to include more than 6 real world problems (problem 43 to 48 newly added), including Colebrook-White, Van der Waals, Proxy for Financial Option Volatility, Beam Deflection, Chemical Equilibrium, orbital mechanics (Kepler problem, page 19) and Wilkinson's polynomial (hard engineering case, page 20). The Abstract/Conclusion (pages 1, 30) claim ``practical root-finding performance across\ldots ill-conditioned equations'' and ``immediately suitable for\ldots real-world applications.'' Future work mentions large-scale engineering (page 30). The benchmark is still mostly synthetic (48 functions, page 15).
\item In-text citations are removed from all parts expcept related work section. The full reference list has been added (pages 31--32), with only 9 entries including [1] Burden et al., [2] Sabharwal, [3] Badr et al., up to [9] Dekker. The GitHub repository is mentioned (page 31) for code.
\item Results use averages over 10,000 runs (page 14), with wins (e.g., 47.9\% by NFE, page 1) and tables (pages 21-26). No hypothesis testing, CI, or std dev shown. Conclusions substantiated by multi-metric comparisons (Iteration count, NFE, FLOPs, CPU), with test hardware and software environment details.
\end{enumerate}

\textbf{Reviewer 2}
\begin{enumerate}
\item Changed to adaptive: $\delta = 10^{-8} \times \max(1.0, |x_{FP}|) + 10^{-15}$ (pages 4, 7, 12-13). Discussed in finite-precision safeguards (page 12), balancing truncation/round-off errors.
\item Used lambdify to NumPy for fair numerical execution (page 13). Reports both CPU (e.g., 1.91x over SciPy Brent, page 1) and NFE (primary metric, pages 22, 26). All in same environment.
\item Related Work (page 2) and Theoretical Analysis (page 11) clarify Trisection's 1/3 reduction at 2 evals cost vs. Bisection's 1/2 at 1 eval. Emphasizes "better bracket placements for subsequent interpolation" implicitly via multi-phase design.
\item Pseudocode updates interval based on sign (e.g., page 6), but no explicit $a < x_{FP} < b$ check (interpolation should ensure it as regula falsi is closed bracting method, with $\epsilon$ guard). Modified secant has strict checks (page 7). Finite-precision section (page 12) adaptive epsilon adds safeguards.
\end{enumerate}

Regards,\
Abdelrahman Ellithy
\end{document}